\documentclass[11pt]{article}

\usepackage{amsmath, amsfonts, amssymb, amsthm, graphicx, mathtools,dsfont}
\usepackage[blocks, affil-it]{authblk}

\usepackage[numbers, square]{natbib}
\usepackage[CJKbookmarks=true,
            bookmarksnumbered=true,
            bookmarksopen=true,
            colorlinks=true,
            citecolor=red,
            linkcolor=blue,
            anchorcolor=red,
            urlcolor=blue]{hyperref}

\usepackage{xcolor}
\newcommand{\red}{\color{red}}
\newcommand{\blue}{\color{blue}}

\usepackage[letterpaper, left=1.2truein, right=1.2truein, top=1.2truein, bottom=1.2truein]{geometry}
\usepackage[ruled, vlined, lined, commentsnumbered]{algorithm2e}

\usepackage{prettyref, soul}
\usepackage{enumitem}
\newcommand{\indep}{\mathrel{\perp\mkern-10mu\perp}}
\newtheorem{lemma}{Lemma}
\newtheorem{proposition}{Proposition}
\newtheorem{thm}{Theorem}
\newtheorem{definition}{Definition}
\newtheorem{corollary}{Corollary}
\newtheorem{remark}{Remark}

\newrefformat{eq}{(\ref{#1})}
\newrefformat{chap}{Chapter~\ref{#1}}
\newrefformat{sec}{Section~\ref{#1}}
\newrefformat{algo}{Algorithm~\ref{#1}}
\newrefformat{fig}{Fig.~\ref{#1}}
\newrefformat{tab}{Table~\ref{#1}}
\newrefformat{rmk}{Remark~\ref{#1}}
\newrefformat{clm}{Claim~\ref{#1}}
\newrefformat{def}{Definition~\ref{#1}}
\newrefformat{cor}{Corollary~\ref{#1}}
\newrefformat{lmm}{Lemma~\ref{#1}}
\newrefformat{lemma}{Lemma~\ref{#1}}
\newrefformat{prop}{Proposition~\ref{#1}}
\newrefformat{app}{Appendix~\ref{#1}}
\newrefformat{ex}{Example~\ref{#1}}
\newrefformat{exer}{Exercise~\ref{#1}}
\newrefformat{soln}{Solution~\ref{#1}}
\newrefformat{cond}{Condition~\ref{#1}}

\def\pr{\textsf{P}} 
\def\ep{\textsf{E}} 
\def\Cov{\textsf{Cov}} 
\def\Var{\operatorname{Var}} 
\def\Cal#1{{\mathcal #1}}
\def\bk#1{{\mathbf #1}}
\def\bkg#1{\mbox{\boldmath{$#1$}}} 
\def\smallbkg#1{\mbox{\scriptsize \boldmath{$#1$}}} 

\def\mb{\mathbf}
\def\mr{\mathrm}
\def\dsum{\displaystyle\sum}
\def\dint{\displaystyle\int}
\def\dfrac{\displaystyle\frac}
\def\bd{\boldsymbol}

\DeclarePairedDelimiter{\ceil}{\lceil}{\rceil}
\newcommand{\bbR}{\mathbb{R}}
\newcommand{\argmin}{\mathop{\rm argmin}}
\newcommand{\argmax}{\mathop{\rm argmax}}
\newcommand{\lnorm}[2]{\left\|{#1} \right\|_{{#2}}}
\newcommand{\indc}[1]{{\mathds{1}_{\left\{{#1}\right\}}}}
\newcommand{\norm}[1]{\|{#1} \|}
\newcommand{\hf}{{1/2}}
\newcommand{\wh}{\widehat}
\newcommand{\wt}{\widetilde}
\newcommand{\Norm}[1]{\left\|{#1}\right\|}
\newcommand{\Fnorm}[1]{\lnorm{#1}{\rm F}}
\newcommand{\opnorm}[1]{\lnorm{#1}{\rm op}}
\newcommand{\nnorm}[1]{\lnorm{#1}{\rm *}}
\newcommand{\Prob}{\mathbb{P}}
\newcommand{\Expect}{\mathbb{E}}
\newcommand{\Span}{{\rm span}}
\newcommand{\rank}{\mathop{\sf rank}}
\newcommand{\Tr}{\mathop{\sf Tr}}
\newcommand{\diag}{\mathop{\text{diag}}}
\newcommand{\supp}{{\rm supp}}
\newcommand{\iprod}[2]{\left \langle #1, #2 \right\rangle}
\newcommand{\pth}[1]{\left( #1 \right)}
\newcommand{\qth}[1]{\left[ #1 \right]}
\newcommand{\sth}[1]{\left\{ #1 \right\}}
\newcommand{\calL}{\mathcal{L}}
\newcommand{\iid}{\overset{\text{iid}}{\sim}}
\newcommand{\nb}[1]{\texttt{\blue[#1]}}
\newcommand{\E}[1]{\mathbb E_{#1}}
\newcommand{\abs}[1]{\left|#1\right|}
\newcommand{\solution}{\par\vspace{10pt}\noindent\textsc{Solution.}\ }

\newcommand{\lecture}[4]{
   \pagestyle{myheadings}
   \thispagestyle{plain}
   \newpage
   \setcounter{page}{1}
   \noindent
   \begin{center}
   \framebox[\textwidth]{
      \begin{minipage}{0.96\textwidth}
        \vspace{2mm}
        \textbf{30815: Asymptotic Theory for Time Series Analysis} \hfill \textbf{Fall 2026} \\
        \textbf{\phantom{aaaaaa} and Stochastic Approximation Algorithms}
        \vspace{4mm}
        \begin{center}
            {\Large Lecture Notes #1 #2}
        \end{center}
        \vspace{2mm}
        \textit{Lecturer: #3} \hfill \textit{Student: #4}
        \vspace{2mm}
      \end{minipage}
   }
   \end{center}
   \markboth{Lecture Notes #1}{Lecture Notes #1}
   \vspace*{4mm}
}

\begin{document}

\lecture{11/09}{}{Wei Biao Wu}{Giorgio Micaletto}
\noindent
We turn to covariance matrices of i.i.d.\ high-dimensional vectors. When the dimension is comparable to the sample size, the sample covariance matrix is inconsistent in operator norm: its spectrum spreads out according to the Marchenko--Pastur law, and its largest eigenvalue fluctuates on the Tracy--Widom scale. In general the error is governed by the effective rank of $\Sigma$, and consistency can be restored by regularization: banding \citep{BL08a} or thresholding \citep{BL08b}.

\paragraph{Setting and notation.} The observations $X_1,\dots,X_n\in\bbR^p$ are i.i.d.\ with $\Expect X_i=0$, covariance $\Sigma=\Cov(X_i)=(\sigma_{jk})$ and precision matrix $\Omega=\Sigma^{-1}$. Write $X_i=(X_{i1},\dots,X_{ip})^\top$. With the mean known to be $0$, the sample covariance matrix is
\[
\wh\Sigma_n=\frac1n\sum_{i=1}^nX_iX_i^\top=(\wh\sigma_{jk})_{j,k=1}^p,\qquad \wh\sigma_{jk}=\frac1n\sum_{i=1}^nX_{ij}X_{ik}.
\]
For a symmetric matrix $M$, $\rho(M)=\max_l|\lambda_l(M)|=\opnorm M$ denotes its spectral radius. We repeatedly use
\begin{equation}\label{eq:rowsum}
\rho(M)\le\max_k\sum_{j=1}^p|m_{jk}|\qquad(M=M^\top),
\end{equation}
which holds because every eigenvalue lies in a Gershgorin disc.

%--------------------------------------------------------------------
\paragraph{The sample covariance matrix in high dimensions:}

\begin{enumerate}[label=(\roman*)]
\item If $p>n$, then $\rank(\wh\Sigma_n)\le n<p$, so $\wh\Sigma_n$ is singular and $\wh\Omega_n=\wh\Sigma_n^{-1}$ does not exist.
\item If $p$ is fixed, $\wh\Sigma_n\to\Sigma$ almost surely by the law of large numbers, also in operator norm.
\item If $p$ is large, this breaks down even in the simplest case. Let $p=n$ and $X_{ij}$ i.i.d.\ $N(0,1)$, so that $\Sigma$ is the identity. Then $\rho(\wh\Sigma_n)\to4$ almost surely \citep{Gem80,YBK88}, while $\rho(\Sigma)=1$. Since also $\lambda_{\min}(\wh\Sigma_n)\to0$, the eigenvalues of $\wh\Sigma_n-\Sigma$ fill out $[-1,3]$ and $\rho(\wh\Sigma_n-\Sigma)\to3$. Moreover $\wh\Sigma_n^{-1}$, although it exists for $p=n$, has $\rho(\wh\Sigma_n^{-1})\to\infty$. For $p/n\to y\in(0,\infty)$ the limit of $\rho(\wh\Sigma_n)$ is $(1+\sqrt y)^2$.
\end{enumerate}

%--------------------------------------------------------------------
\paragraph{Marchenko--Pastur.}

Let $\wh\lambda_1\ge\dots\ge\wh\lambda_p$ be the eigenvalues of $\wh\Sigma_n$ and define their empirical c.d.f.\ (the empirical spectral distribution, ESD)
\[
\wh F_n(t)=\frac1p\sum_{l=1}^p\mathbf 1\{\wh\lambda_l\le t\}.
\]

\begin{thm}[Marchenko--Pastur]\label{thm:mp}
Let the $X_{ij}$ be i.i.d.\ with mean $0$ and variance $1$, and $p/n\to y\in(0,\infty)$. Then $\wh F_n\Rightarrow F_y$ almost surely, where $F_y$ has density
\[
f_y(t)=\frac{\sqrt{(b-t)(t-a)}}{2\pi y\,t},\qquad a\le t\le b,\qquad a=(1-\sqrt y)^2,\ b=(1+\sqrt y)^2,
\]
plus an atom of mass $1-1/y$ at $0$ when $y>1$. In particular, for $p=n$ ($y=1$),
\[
f(t)=\frac{\sqrt{t(4-t)}}{2\pi t}=\frac{1}{2\pi}\,t^{-1/2}(4-t)^{1/2},\qquad0<t\le4,
\]
i.e.\ $F_1$ is the law of $4U$ with $U\sim\mathrm{Beta}(\frac12,\frac32)$.
\end{thm}

See \citet{MP67} and \citet[Ch.~3]{BS10}. On the board the density read $\sqrt{t(4-t)}/(\pi t)$; the constant must be $2\pi$ for it to integrate to $1$, since $\int_0^4\sqrt{(4-t)/t}\,dt=2\pi$ (substitute $t=4\sin^2\varphi$). The Beta form follows from $B(\frac12,\frac32)=\pi/2$. Although the variables are independent with variance $1$, the eigenvalues of $\wh\Sigma_n$ spread over the whole interval $[0,4]$ (Figure~\ref{fig:mp}).

\begin{figure}[ht]
\centering
\setlength{\unitlength}{1mm}
\begin{picture}(90,48)
\put(10,5){\vector(1,0){72}}
\put(10,5){\vector(0,1){40}}
\put(83,5){\makebox(0,0)[l]{$t$}}
\put(10,46){\makebox(0,0)[b]{$f(t)$}}
\qbezier(10.75,40.37)(11.12,35.11)(11.50,29.85)
\qbezier(11.50,29.85)(11.88,27.50)(12.25,25.16)
\qbezier(12.25,25.16)(12.62,23.75)(13.00,22.34)
\qbezier(13.00,22.34)(13.38,21.38)(13.75,20.41)
\qbezier(13.75,20.41)(14.12,19.69)(14.50,18.97)
\qbezier(14.50,18.97)(14.88,18.41)(15.25,17.85)
\qbezier(15.25,17.85)(15.62,17.39)(16.00,16.94)
\qbezier(16.00,16.94)(16.38,16.56)(16.75,16.18)
\qbezier(16.75,16.18)(17.12,15.85)(17.50,15.53)
\qbezier(17.50,15.53)(18.25,15.00)(19.00,14.47)
\qbezier(19.00,14.47)(20.50,13.71)(22.00,12.96)
\qbezier(22.00,12.96)(23.50,12.42)(25.00,11.89)
\qbezier(25.00,11.89)(26.50,11.48)(28.00,11.08)
\qbezier(28.00,11.08)(29.50,10.75)(31.00,10.42)
\qbezier(31.00,10.42)(32.50,10.15)(34.00,9.87)
\qbezier(34.00,9.87)(35.50,9.64)(37.00,9.40)
\qbezier(37.00,9.40)(38.50,9.19)(40.00,8.98)
\qbezier(40.00,8.98)(41.50,8.79)(43.00,8.60)
\qbezier(43.00,8.60)(44.50,8.42)(46.00,8.25)
\qbezier(46.00,8.25)(47.50,8.08)(49.00,7.92)
\qbezier(49.00,7.92)(50.50,7.76)(52.00,7.60)
\qbezier(52.00,7.60)(53.50,7.45)(55.00,7.30)
\qbezier(55.00,7.30)(56.50,7.14)(58.00,6.99)
\qbezier(58.00,6.99)(59.50,6.83)(61.00,6.67)
\qbezier(61.00,6.67)(62.50,6.50)(64.00,6.33)
\qbezier(64.00,6.33)(65.50,6.12)(67.00,5.91)
\qbezier(67.00,5.91)(68.50,5.46)(70.00,5.00)
\put(70,4){\line(0,1){2}}
\put(70,2.5){\makebox(0,0)[t]{\small $4$}}
\put(10,2.5){\makebox(0,0)[t]{\small $0$}}
\put(25,4){\line(0,1){2}}
\put(25,2.5){\makebox(0,0)[t]{\small $1$}}
\put(32,22){\makebox(0,0)[l]{\small $f(t)=\dfrac{\sqrt{t(4-t)}}{2\pi t}$}}
\end{picture}
\caption{Marchenko--Pastur density for $y=1$: it blows up like $t^{-1/2}$ at $0$ and vanishes at the right edge $4$.}
\label{fig:mp}
\end{figure}

%--------------------------------------------------------------------
\paragraph{The moment method.}

The moments of the ESD are normalized traces:
\[
\int_0^\infty t^k\,d\wh F_n(t)=\frac1p\sum_{l=1}^p\wh\lambda_l^k=\frac1p\Tr\big(\wh\Sigma_n^k\big),\qquad
\beta_k:=\int_0^\infty t^k\,dF(t)=\int_0^\infty t^kf(t)\,dt .
\]

\begin{lemma}[Moment convergence theorem]\label{lmm:moments}
Suppose $F$ is determined by its moments, which holds under Carleman's condition $\sum_{k=1}^\infty\beta_{2k}^{-1/(2k)}=\infty$, in particular if $F$ has compact support. If $\int t^k\,d\wh F_n(t)\to\beta_k$ in probability for every $k=1,2,\dots$, then $\wh F_n\Rightarrow F$ in probability.
\end{lemma}

This is the method of moments \citep[Thm.~30.2]{Bil95}, applied along subsequences. On the board Carleman's condition appeared as $\sum_k\beta_{2k}^{-1/(2k)}<\infty$; the condition for determinacy is divergence of the series. For the Marchenko--Pastur law with $y=1$ it is automatic: the support is $[0,4]$, so $\beta_{2k}^{1/(2k)}\le4$ and every term of the series is at least $1/4$.

The first two moments answer the question marks on the board.

\begin{proposition}\label{prop:twomoments}
Let the $X_{ij}$ be i.i.d.\ $N(0,1)$ and $p/n\to y$. Then
\[
\Expect\Big[\frac1p\Tr\wh\Sigma_n\Big]=1,\qquad
\Expect\Big[\frac1p\Tr\wh\Sigma_n^2\Big]=1+\frac pn+\frac1n\longrightarrow1+y,
\]
and both variances tend to $0$. For $y=1$ the limits $1$ and $2$ are the first two moments of $4\,\mathrm{Beta}(\frac12,\frac32)$.
\end{proposition}

\begin{proof}
First, $\frac1p\Tr\wh\Sigma_n=\frac1{np}\sum_{i,j}X_{ij}^2$ has mean $1$ and variance $2/(np)$. Next,
\[
\Tr\wh\Sigma_n^2=\frac1{n^2}\sum_{j,k=1}^p\Big(\sum_{i=1}^nX_{ij}X_{ik}\Big)^2 .
\]
For $j\ne k$ the inner square has mean $n$; for $j=k$ it has mean $\Expect(\sum_iX_{ij}^2)^2=3n+n(n-1)=n^2+2n$. Hence
\[
\Expect\Tr\wh\Sigma_n^2=n^{-2}\big(p(p-1)n+p(n^2+2n)\big)=\frac{p^2}n+p+\frac pn .
\]
The variance bound is a similar but longer computation. If $U\sim\mathrm{Beta}(a,b)$ then $\Expect U^k=\prod_{r=0}^{k-1}\frac{a+r}{a+b+r}$; with $a=\frac12$, $b=\frac32$ this gives $\Expect(4U)=1$ and $\Expect(4U)^2=2$.
\end{proof}

\begin{remark}[All moments]
In general $\frac1p\Tr\wh\Sigma_n^k\to\beta_k(y)=\sum_{r=0}^{k-1}\frac{y^r}{r+1}\binom kr\binom{k-1}r$ in probability \citep[Sec.~3.1]{BS10}. Expanding
\[
\Tr\big(\wh\Sigma_n^k\big)=n^{-k}\sum X_{i_1j_1}X_{i_1j_2}X_{i_2j_2}X_{i_2j_3}\cdots X_{i_kj_k}X_{i_kj_1}
\]
as a sum over closed walks $j_1\to i_1\to j_2\to\dots\to i_k\to j_1$ on a bipartite graph, the leading contributions come from walks that traverse each edge exactly twice and whose edges form a tree; counting them gives $\beta_k(y)$. For $y=1$ these are the Catalan numbers $\beta_k=\frac1{k+1}\binom{2k}k=1,2,5,14,42,\dots$, which are exactly the moments of $4U$: $\Expect(4U)^k=4^k\prod_{r=0}^{k-1}\frac{1/2+r}{2+r}=\frac{(2k)!}{k!\,(k+1)!}$. By Lemma~\ref{lmm:moments}, $\wh F_n\Rightarrow F_1$. The proof via moments needs all moments of $X_{ij}$; finite variance suffices after truncation.
\end{remark}

\begin{remark}[The largest eigenvalue via moments]
Convergence of the ESD only gives $\liminf\wh\lambda_1\ge4$. For an upper bound, note
\[
\wh\lambda_1^k\le\sum_{l=1}^p\wh\lambda_l^k=\Tr\big(\wh\Sigma_n^k\big)\le p\,\wh\lambda_1^k,
\qquad\text{so}\qquad
\wh\lambda_1\le\big(\Tr\wh\Sigma_n^k\big)^{1/k}\le p^{1/k}\,\wh\lambda_1 .
\]
Once $k\gg\log p$, $p^{1/k}\to1$ and $\sum_l\wh\lambda_l^k\approx\wh\lambda_1^k$. Bounding $\Expect\Tr(\wh\Sigma_n^{k_n})$ with $k=k_n\to\infty$ is how \citet{Gem80} proved $\wh\lambda_1\to(1+\sqrt y)^2$.
\end{remark}

%--------------------------------------------------------------------
\paragraph{Tracy--Widom.}

\begin{thm}[\citet{Joh01}]\label{thm:tw}
Let the $X_{ij}$ be i.i.d.\ $N(0,1)$ and $p=n$. Then
\[
n^{2/3}\big(\wh\lambda_1-4\big)\Rightarrow2^{4/3}\,W_1,\qquad W_1\sim F_1\ \text{(Tracy--Widom law of order 1)}.
\]
More generally, with $\mu_{np}=(\sqrt{n-1}+\sqrt p)^2$ and $\sigma_{np}=(\sqrt{n-1}+\sqrt p)\big(\frac1{\sqrt{n-1}}+\frac1{\sqrt p}\big)^{1/3}$, $(n\wh\lambda_1-\mu_{np})/\sigma_{np}\Rightarrow F_1$ when $n/p$ converges to a positive constant (see \citet{Joh01} for the precise conditions). For $p=n$, $\mu_{np}/n=4+O(n^{-1})$ and $\sigma_{np}\sim2^{4/3}n^{1/3}$, which gives the display.
\end{thm}

The Tracy--Widom distribution functions are \citep{TW94,TW96}
\[
F_2(x)=\exp\Big(-\int_x^\infty(t-x)\,q(t)^2\,dt\Big),\qquad
F_1(x)=\exp\Big(-\frac12\int_x^\infty\big[q(t)+(t-x)\,q(t)^2\big]\,dt\Big),
\]
where $q$ solves the Painlevé II equation $q''=tq+2q^3$ with $q(t)\sim\mathrm{Ai}(t)$ as $t\to\infty$ (the Hastings--McLeod solution). On the board, $e^{\int_x^\infty(t-x)^2q^2(t)\,dt}$ was labeled as the pdf. The correct formula is the c.d.f.\ $F_2$ above: a minus sign in the exponent and the factor $(t-x)$ without the square. $F_2$ is the limit for complex Gaussian data \citep{Joh00}; for real data, as here, the limit is $F_1$. Note the scale: the largest eigenvalue fluctuates at order $n^{-2/3}$, much smaller than the usual $n^{-1/2}$.

%--------------------------------------------------------------------
\paragraph{Effective rank.}
For general $\Sigma$ the coordinates $X_{i1},\dots,X_{ip}$ are dependent, and the relevant measure of dimension is the effective rank
\[
r(\Sigma)=\frac{\Tr(\Sigma)}{\rho(\Sigma)}\in[1,p].
\]

\begin{thm}[\citet{KL17}]\label{thm:kl}
If the $X_i$ are centered Gaussian vectors, then, with absolute constants in $\asymp$,
\[
\Expect\,\rho\big(\wh\Sigma_n-\Sigma\big)\asymp\rho(\Sigma)\Big(\sqrt{\frac{r(\Sigma)}n}\ \vee\ \frac{r(\Sigma)}n\Big).
\]
\end{thm}

On the board this appeared as $\rho(\wh\Sigma_n-\Sigma)\asymp\rho(\Sigma)\sqrt{r(\Sigma)/n}$, which is the case $r(\Sigma)\le n$. The relative error $\rho(\wh\Sigma_n-\Sigma)/\rho(\Sigma)$ therefore tends to $0$ if and only if $r(\Sigma)=o(n)$: what matters is the effective rank, not the ambient dimension $p$. Two examples:
\begin{itemize}
\item $\Sigma=\bk I_p$: $r(\Sigma)=p$, so the error is $\asymp\sqrt{p/n}\vee p/n$, which does not vanish when $p\asymp n$, consistent with the large $p$ regime.
\item $\Sigma=\mathbf 1\mathbf 1^\top$ (all entries $1$): $\Tr(\Sigma)=p=\rho(\Sigma)$, so $r(\Sigma)=1$ and the relative error is $\asymp n^{-1/2}$ for every $p$. Indeed $X_i=Z_i\mathbf 1$ with $Z_i\sim N(0,1)$, and $\wh\Sigma_n-\Sigma=(\frac1n\sum_iZ_i^2-1)\,\mathbf 1\mathbf 1^\top$.
\end{itemize}

%--------------------------------------------------------------------
\paragraph{Banding.}
Suppose the variables have a natural order (time, space, wavelength) and $\Sigma$ can be approximated by banded matrices. For $\alpha>0$ consider the bandable classes
\begin{equation}\label{eq:bandable}
\max_k\sum_{j:\,|j-k|>B}|\sigma_{jk}|\le C\,B^{-\alpha}\quad\text{for all }B\ge1
\qquad\text{or}\qquad
\max_k\sum_{j:\,|j-k|>B}|\sigma_{jk}|\le C_1e^{-C_2B}.
\end{equation}
Examples:
\begin{itemize}
\item Covariances $\sigma_{jk}=f(|j-k|)$ with $|f(d)|\le Cd^{-1-\alpha}$ fall in the polynomial class.
\item Exponentially decaying $f$ fall in the exponential class. This includes the AR(1) covariance $\sigma_{jk}\propto\phi^{|j-k|}$, $|\phi|<1$, and the Matérn covariance functions of spatial statistics \citep{Ste99}.
\end{itemize}

For AR-type dependence it is the precision matrix $\Omega$, rather than $\Sigma$, that is banded (tridiagonal for an AR(1)). One then bands $\Omega$ through its modified Cholesky decomposition, which amounts to fitting AR regressions \citep{WP03,BL08a}.


\medskip
The banded estimator of \citet{BL08a} is
\[
\wt\Sigma_{n,B}=\big(\wh\sigma_{jk}\,\mathbf 1\{|j-k|\le B\}\big)_{j,k=1}^p,
\]
the Schur (entrywise) product of $\wh\Sigma_n$ with the banding mask.

\begin{lemma}[Bernstein's inequality]\label{lmm:bernstein}
Let $Z_1,\dots,Z_n$ be i.i.d.\ with $\Expect Z_1=0$ and $\Expect|Z_1|^m\le\frac{m!}2K^{m-2}\Expect Z_1^2$ for $m\ge2$ (e.g.\ $Z_1$ sub-exponential). Then for $u>0$,
\[
\Prob\Big(\Big|\sum_{i=1}^nZ_i\Big|\ge u\Big)\le2\exp\Big(-\frac{u^2}{C_1u+C_2\,n\Expect Z_1^2}\Big),\qquad C_1=2K,\ C_2=2 .
\]
\end{lemma}

\begin{thm}[\citet{BL08a}]\label{thm:band}
Let $X_i\sim N(0,\Sigma)$ with $\max_j\sigma_{jj}\le M$ and $\Sigma$ in the polynomial class \eqref{eq:bandable}, and let $p=p_n\to\infty$ with $\log p/n\to0$. Then
\[
\rho\big(\wt\Sigma_{n,B}-\Sigma\big)=O_\Prob\Big(B\sqrt{\frac{\log p}n}+B^{-\alpha}\Big).
\]
In particular $\rho(\wt\Sigma_{n,B}-\Sigma)\to0$ in probability whenever $B=B_n\to\infty$ and $B_n\sqrt{\log p/n}\to0$. The choice $B\asymp(n/\log p)^{1/(2(\alpha+1))}$ gives the rate $(\log p/n)^{\alpha/(2(\alpha+1))}$.
\end{thm}

\begin{proof}
Split into a variance part and a bias part:
\[
\rho\big(\wt\Sigma_{n,B}-\Sigma\big)\le\rho\big(\wt\Sigma_{n,B}-\Expect\wt\Sigma_{n,B}\big)+\rho\big(\Expect\wt\Sigma_{n,B}-\Sigma\big).
\]
\emph{Bias.} Since $\Expect\wh\sigma_{jk}=\sigma_{jk}$, the matrix $\Expect\wt\Sigma_{n,B}-\Sigma$ is $0$ inside the band and $-\sigma_{jk}$ outside. By \eqref{eq:rowsum} and \eqref{eq:bandable},
\[
\rho\big(\Expect\wt\Sigma_{n,B}-\Sigma\big)\le\max_k\sum_{j:\,|j-k|>B}|\sigma_{jk}|\le CB^{-\alpha}.
\]
\emph{Variance.} By \eqref{eq:rowsum},
\[
\rho\big(\wt\Sigma_{n,B}-\Expect\wt\Sigma_{n,B}\big)\le\max_k\sum_{j:\,|j-k|\le B}|\wh\sigma_{jk}-\sigma_{jk}|\le(2B+1)\max_{j,k}|\wh\sigma_{jk}-\sigma_{jk}| .
\]
We need $y$ such that $\Prob\big(\max_k\sum_{|j-k|\le B}|\wh\sigma_{jk}-\sigma_{jk}|\ge y\big)\to0$. A union bound over the $p$ rows, and then over the at most $2B+1\le p$ entries in each row, gives
\begin{align*}
\Prob\Big(\max_k\sum_{|j-k|\le B}|\wh\sigma_{jk}-\sigma_{jk}|\ge y\Big)
&\le p\max_k\Prob\Big(\sum_{|j-k|\le B}|\wh\sigma_{jk}-\sigma_{jk}|\ge y\Big)\\
&\le p^2\max_{j,k}\Prob\Big(|\wh\sigma_{jk}-\sigma_{jk}|\ge\frac y{2B+1}\Big).
\end{align*}
Now $n(\wh\sigma_{jk}-\sigma_{jk})=\sum_{i=1}^nZ_i$ with $Z_i=X_{ij}X_{ik}-\sigma_{jk}$ i.i.d., centered and sub-exponential, with $\Expect Z_1^2=\sigma_{jj}\sigma_{kk}+\sigma_{jk}^2\le2M^2$ and a constant $K$ depending only on $M$. With $u=ny/(2B+1)$, Lemma~\ref{lmm:bernstein} gives the bound
\[
2p^2\exp\Big(-\frac{u^2}{C_1u+C_2n\Expect Z_1^2}\Big).
\]
Choose $u$ so that this tends to $0$. If $C_1u$ dominates the denominator, the exponent is at least $u/(2C_1)$, so we need $p^2e^{-u/(2C_1)}\to0$, i.e.\ $u\ge C\log p$. If $C_2n\Expect Z_1^2$ dominates, the exponent is at least $u^2/(2C_2n\Expect Z_1^2)$, so we need $u\ge C\sqrt{n\log p}$. Taking
\[
u=C\big(\sqrt{n\log p}+\log p\big)
\]
with a big constant $C$, both exponents are at least $3\log p$ and the bound is at most $2/p\to0$. The corresponding threshold is
\[
y=\frac{(2B+1)u}n\le C'B\,\frac{\sqrt{n\log p}+\log p}n=C'B\Big(\sqrt{\frac{\log p}n}+\frac{\log p}n\Big),
\]
which is $O(B\sqrt{\log p/n})$ because $\log p/n\to0$. Combining the two parts proves the bound, and balancing $B\sqrt{\log p/n}$ against $B^{-\alpha}$ gives the optimal $B$.
\end{proof}

\begin{remark}
The dimension enters only through $\log p$, the price of the union bound over $p^2$ entries, so $p$ may grow exponentially in $n$. Compare the time-series case of the last lecture: there a single realization and the Toeplitz structure lead to rates $B/\sqrt n$, or $\sqrt{B\log B/n}$; here we have $n$ replicates but no structure across entries. \citet{BL08a} also prove the rate for $\wt\Sigma_{n,B}^{-1}$ when $\Sigma$ is well conditioned, and extend the result to sub-Gaussian tails.
\end{remark}

%--------------------------------------------------------------------
\paragraph{Nagaev's inequality.}

Bernstein's inequality requires exponential moments. If $Z_i$ only has a finite $q$-th moment, $q>2$, Nagaev's inequality \citep{Nag79} (Lemma~1.2 in the notes of 11/02) gives, for i.i.d.\ centered $Z_i$,
\[
\Prob\Big(\Big|\sum_{i=1}^nZ_i\Big|\ge y\Big)\le C_1\,\frac{n\,\Expect|Z_1|^q}{y^q}+C_2\exp\Big(-C_3\,\frac{y^2}{n\,\Expect Z_1^2}\Big).
\]
On the board the factor $n$ in the polynomial term was missing. The exponential term reproduces the Gaussian-type bound used above; the polynomial term is new.

%--------------------------------------------------------------------
\paragraph{Thresholding.}

Banding uses the order of the variables. When there is no natural order but $\Sigma$ is sparse, one keeps only the large entries:
\[
\text{banded: }\big(\wh\sigma_{jk}\,\mathbf 1\{|j-k|\le B\}\big),\qquad
\text{thresholded: }T_u(\wh\Sigma_n)=\big(\wh\sigma_{jk}\,\mathbf 1\{|\wh\sigma_{jk}|\ge u\}\big).
\]
Thresholding is invariant under permutations of the variables. By the argument in the proof of Theorem~\ref{thm:band}, $\max_{j,k}|\wh\sigma_{jk}-\sigma_{jk}|=O_\Prob(\sqrt{\log p/n})$, so a threshold $u\asymp\sqrt{\log p/n}$ removes the entries that are pure noise.


\begin{thm}[\citet{BL08b}]
Let $X_i\sim N(0,\Sigma)$, $0\le q<1$, and consider the class $\mathcal U(q,c_0(p),M)=\{\Sigma:\sigma_{jj}\le M,\ \max_j\sum_k|\sigma_{jk}|^q\le c_0(p)\}$. If $u=M'\sqrt{\log p/n}$ with $M'$ large enough and $\log p/n\to0$, then uniformly on this class
\[
\rho\big(T_u(\wh\Sigma_n)-\Sigma\big)=O_\Prob\Big(c_0(p)\Big(\frac{\log p}n\Big)^{(1-q)/2}\Big).
\]
\end{thm}

For $q=0$, $c_0(p)$ bounds the number of nonzero entries per row. See also \citet{EK08}.

%--------------------------------------------------------------------
\begin{thebibliography}{99}

\bibitem[Bai and Silverstein(2010)]{BS10}
Z.~Bai and J.~W. Silverstein.
\newblock \emph{Spectral Analysis of Large Dimensional Random Matrices}.
\newblock Springer, New York, 2nd edition, 2010.

\bibitem[Bickel and Levina(2008a)]{BL08a}
P.~J. Bickel and E.~Levina.
\newblock Regularized estimation of large covariance matrices.
\newblock \emph{Ann. Statist.}, 36:199--227, 2008a.

\bibitem[Bickel and Levina(2008b)]{BL08b}
P.~J. Bickel and E.~Levina.
\newblock Covariance regularization by thresholding.
\newblock \emph{Ann. Statist.}, 36:2577--2604, 2008b.

\bibitem[Billingsley(1995)]{Bil95}
P.~Billingsley.
\newblock \emph{Probability and Measure}.
\newblock Wiley, New York, 3rd edition, 1995.

\bibitem[El~Karoui(2008)]{EK08}
N.~El~Karoui.
\newblock Operator norm consistent estimation of large-dimensional sparse covariance matrices.
\newblock \emph{Ann. Statist.}, 36:2717--2756, 2008.

\bibitem[Geman(1980)]{Gem80}
S.~Geman.
\newblock A limit theorem for the norm of random matrices.
\newblock \emph{Ann. Probab.}, 8:252--261, 1980.

\bibitem[Johansson(2000)]{Joh00}
K.~Johansson.
\newblock Shape fluctuations and random matrices.
\newblock \emph{Comm. Math. Phys.}, 209:437--476, 2000.

\bibitem[Johnstone(2001)]{Joh01}
I.~M. Johnstone.
\newblock On the distribution of the largest eigenvalue in principal components analysis.
\newblock \emph{Ann. Statist.}, 29:295--327, 2001.

\bibitem[Koltchinskii and Lounici(2017)]{KL17}
V.~Koltchinskii and K.~Lounici.
\newblock Concentration inequalities and moment bounds for sample covariance operators.
\newblock \emph{Bernoulli}, 23:110--133, 2017.

\bibitem[Mar\v{c}enko and Pastur(1967)]{MP67}
V.~A. Mar\v{c}enko and L.~A. Pastur.
\newblock Distribution of eigenvalues for some sets of random matrices.
\newblock \emph{Math. USSR-Sb.}, 1:457--483, 1967.

\bibitem[Nagaev(1979)]{Nag79}
S.~V. Nagaev.
\newblock Large deviations of sums of independent random variables.
\newblock \emph{Ann. Probab.}, 7:745--789, 1979.

\bibitem[Stein(1999)]{Ste99}
M.~L. Stein.
\newblock \emph{Interpolation of Spatial Data: Some Theory for Kriging}.
\newblock Springer, New York, 1999.

\bibitem[Tracy and Widom(1994)]{TW94}
C.~A. Tracy and H.~Widom.
\newblock Level-spacing distributions and the {A}iry kernel.
\newblock \emph{Comm. Math. Phys.}, 159:151--174, 1994.

\bibitem[Tracy and Widom(1996)]{TW96}
C.~A. Tracy and H.~Widom.
\newblock On orthogonal and symplectic matrix ensembles.
\newblock \emph{Comm. Math. Phys.}, 177:727--754, 1996.

\bibitem[Wu and Pourahmadi(2003)]{WP03}
W.~B. Wu and M.~Pourahmadi.
\newblock Nonparametric estimation of large covariance matrices of longitudinal data.
\newblock \emph{Biometrika}, 90:831--844, 2003.

\bibitem[Yin et~al.(1988)Yin, Bai, and Krishnaiah]{YBK88}
Y.~Q. Yin, Z.~D. Bai, and P.~R. Krishnaiah.
\newblock On the limit of the largest eigenvalue of the large dimensional sample covariance matrix.
\newblock \emph{Probab. Theory Related Fields}, 78:509--521, 1988.

\end{thebibliography}

\end{document}